\documentclass{article}

\usepackage[utf8]{inputenc}
\usepackage[T1]{fontenc}
\usepackage{amsmath,amssymb,amsthm}
\usepackage{mathtools}
\usepackage{hyperref}
\usepackage{booktabs}

\makeatletter
\newcommand{\citep}[2][]{%
  \begingroup
  \def\@tmpkey{#2}%
  \ifx\@tmpkey\@empty\else
    (\ifx\@empty#1\@empty\else #1, \fi
     \@nameuse{bibkey@\@tmpkey})%
  \fi
  \endgroup
}
\expandafter\def\csname bibkey@lasser2026exo\endcsname{Exogenous~Verification}
\expandafter\def\csname bibkey@lasser2026phase\endcsname{Phase~Redundancy}
\expandafter\def\csname bibkey@lasser2026micro\endcsname{Microfoundation}
\makeatother

\newtheorem{definition}{Definition}
\newtheorem{lemma}{Lemma}
\newtheorem{theorem}{Theorem}
\newtheorem{remark}{Remark}
\newtheorem{assumption}{Assumption}

\newcommand{\Aadv}{A_{\mathrm{adv}}}
\newcommand{\deltaadv}{\delta_{\mathrm{adv}}}
\newcommand{\deltastar}{\delta_{*}}
\newcommand{\rhogap}{\rho_{\mathrm{gap}}}
\newcommand{\Pproxy}{P}
\newcommand{\Ttruth}{T}
\newcommand{\epsgap}{\varepsilon_{\mathrm{gap}}}
\newcommand{\epsnonres}{\varepsilon_{\mathrm{gap}}^{\mathrm{nonres}}}
\newcommand{\epsfloor}{\varepsilon_{\mathrm{floor}}^{\mathrm{res}}}

\newcommand{\Tbeta}{T_{\beta}}
\newcommand{\Ncasc}{N_{\mathrm{cascade}}}
\newcommand{\Nev}{N_{\mathrm{events}}}
\newcommand{\bclk}{\beta_{\mathrm{clk}}}

\newcommand{\bcal}{\beta_{\mathrm{cal}}}
\newcommand{\Bclip}{B}
\newcommand{\Rh}{R_{\mathrm{H}}}
\newcommand{\Kch}{K_{\mathrm{ch}}}

\title{Lemma 6 ($h_{\mathrm{detect}}$ intensivity): Full Proof Draft v6}
\author{Paper 10 working draft for codex re-review}
\date{}

\begin{document}

\maketitle

\section{Goal and changes from v5}

Round~E confirmed Q1 (rigorous Hoeffding inversion) as Sound.  Four
narrow issues remain; this v6 addresses all four plus the
integration-prerequisite expansion.

\textbf{Round E narrow issues.}
\begin{itemize}
\item \textbf{Q2 ($R$ range-width mismatch).}  Lemma~4
  (\texttt{appendices/proofs.tex} §A.4) uses Hoeffding on bounded
  increments with range $R = b - a$ (the interval width).  V5
  clipped to $[-R, R]$ which has width $2R$, making Step~5's
  constant too optimistic by another factor of~4.
\item \textbf{Q3 ((C5.IID) phrasing).}  V5 called raw increments a
  ``martingale-difference sequence'' while also asserting positive
  conditional drift $\geq \deltastar$ — these are inconsistent.
  $\deltastar$ should also be the drift floor for the
  \emph{used/clipped} LLR process, not raw KL from Lemmas~5b/5e
  (clipping can reduce drift).
\item \textbf{Q4 (empirical $\Ncasc$ wrong sign).}  V5 used
  $\Pr[\Ncasc^{\mathrm{est}} \geq \Ncasc^{\mathrm{true}}] \geq 1 -
  \bcal$ — the wrong direction.  For an event-count \emph{floor},
  the operator needs a certified lower bound:
  $\Ncasc^{\mathrm{cert}} \leq \Ncasc^{\mathrm{true}}$ with high
  probability.  Symmetrically, for $\Tbeta$ the operator needs a
  certified upper bound on the true quantile.
\item \textbf{Q5 (bundle into (C5.SPRT)).}  The four sub-clauses
  (HOEFF/MULT/IID/SUPP) read better as a single (C5.SPRT)
  ``SPRT-applicability'' condition with sub-clauses, matching the
  (C9.*) and (C10.*) bundling pattern.  Plus more integration items
  surface (legacy $\beta'$ in Lemma~4/5d statements,
  introduction/theorem statement, §3 invariants connection to
  $I_{11}$, expanded §8 Audit~7).
\end{itemize}

\textbf{Changes from v5 (Round E fixes):}

\begin{itemize}
\item \textbf{Q2 fix ($R$ as Hoeffding range width).}  Renamed
  absolute clip radius to $\Bclip$; defined $\Rh := 2\Bclip$ as
  the Hoeffding range width (matching Lemma~4's convention).
  Step~5 now uses $\Rh$, giving:
  \[
    \Tbeta \;=\; \max\!\left\{\frac{2A}{\deltastar},\;
    \frac{2 \Rh^2 \log(1/\beta)}{\deltastar^2}\right\}
    \;=\; \max\!\left\{\frac{2A}{\deltastar},\;
    \frac{8 \Bclip^2 \log(1/\beta)}{\deltastar^2}\right\}.
  \]
  Plus added one-sentence adversarial-channel-selection clause to
  (C5.MULT): the monitored partition and channel-selection/testing
  policy are fixed before deployment, with familywise
  $\alpha/\beta$ allocation over at most $\Kch$ channels.
\item \textbf{Q3 fix ((C5.IID) cleanup).}  Reformulated: adapted
  increments with conditional mean
  $\mathbb{E}[\ell_n \mid \mathcal{F}_{n-1}] \geq \delta_n \geq
  \deltastar$; centered increments
  $\tilde \ell_n := \ell_n - \mathbb{E}[\ell_n \mid \mathcal{F}_{n-1}]$
  form the martingale-difference sequence used in
  Hoeffding-Azuma.  Also clarified that $\deltastar$ is the
  drift floor for the \emph{used/clipped} LLR process — clipping
  can reduce drift, and operators must verify $\deltastar$ post-
  clipping.
\item \textbf{Q3 fix ((C5.SUPP) cleanup).}  Dropped the numeric
  $\epsilon$ alternative when clipping is mandatory; (C5.SUPP) now
  asserts almost-sure finiteness of $\ell_n$ (which is automatic
  under (C5.HOEFF) clipping).
\item \textbf{Q4 fix (empirical-calibration sign).}  Reformulated
  $\bcal$: in the empirical case, the operator certifies
  conservative bounds on each $\Tbeta$ input ($\Bclip$ upper,
  $\deltastar$ lower, $\alpha/\beta$ allocation, $\Ncasc$ lower).
  Failure probability $\bcal$ is the union of
  $\Pr[\text{any certificate violated}]$.  Total Layer~2 failure
  remains $\beta + \bclk + \bcal$.
\item \textbf{Q5 fix (bundle (C5.SPRT)).}  Bundled HOEFF/MULT/IID/
  SUPP under a single (C5.SPRT) ``SPRT-applicability'' condition
  with sub-clauses (C5.HOEFF), (C5.MULT), (C5.IID), (C5.SUPP),
  matching (C9.*) [BD/CL/WB/TF] and (C10.*) [CN/SU] patterns.
\item \textbf{Integration prerequisites expanded.}  Added §2
  notation table updates, Lemma~4/5d statement updates (not just
  appendix prose), introduction/theorem statement legacy $\beta'$
  drops, §3 invariants/$I_{11}$ connection, expanded Audit~7.
\end{itemize}

\textbf{Carried forward from v5 (unchanged in v6):}
\begin{itemize}
\item Rigorous Hoeffding inversion chain (Step~5 Substeps 5a--5e),
  now with corrected $\Rh = 2\Bclip$.
\item Supremum-form $h_{\mathrm{detect}}$, boundary convention
  $t_0 := t_0^-$, $\rhogap$ rename.
\item Union-class KL floor $\deltastar$ (now applied to clipped
  LLR drift).
\item (C11) framing, (C11.CLK) deterministic floor + $\bclk$
  structure (with v6 sign fix on $\bcal$).
\item Constant taxonomy structure.
\end{itemize}

\subsection*{Earlier-round historical context}

V4 corrected v3's $\kappa$-scaling error and introduced (C5.HOEFF)
+ strengthened (C11.CLK).  Round~D returned ``v5 needed'': the
$1/\deltastar^2$ scaling is conceptually correct, but v4's
specific formula
$\Tbeta = \max\{2A/\deltastar, R^2\log(1/\beta)/(2\deltastar^2)\}$
is not a rigorous inversion of Lemma~4's exact Hoeffding bound:
substituting $t = R^2\log(1/\beta)/(2\deltastar^2)$ does not give
exponent $\geq \log(1/\beta)$ when $A > 0$.

\textbf{Round D significant problem (Q1, factor-of-4 error).}
For $T \geq 2A/\deltastar$, we have $T\deltastar - A \geq
T\deltastar/2$, hence $(T\deltastar - A)^2 \geq T^2\deltastar^2/4$.
Substituting into the Hoeffding bound:
$2(T\deltastar - A)^2/(TR^2) \geq T\deltastar^2/(2R^2)$.
For this to exceed $\log(1/\beta)$, we need $T \geq 2R^2\log(1/\beta)/
\deltastar^2$ — \emph{four times} v4's coefficient.

V5 uses the conservative form
\[
  \Tbeta(\beta, \deltastar) \;:=\; \max\!\left\{
  \frac{2A}{\deltastar},\;
  \frac{2 R^2 \log(1/\beta)}{\deltastar^2}\right\}
\]
which rigorously gives $\Pr[T_{\mathrm{detect}} > \Tbeta] \leq
\beta$.

\textbf{Round D other issues.}
\begin{itemize}
\item Q2: (C5.HOEFF) ``family-wise corrected'' clause too weak —
  Bonferroni controls error probability, not LLR range. Need
  fixed-cardinality monitored partitions or LLR clipping.
\item Q3: (C11.CLK) sound if $\Ncasc \geq \Tbeta$ is hard audit
  constraint; if empirically estimated, calibration uncertainty
  needs separate $\bcal$ failure probability.
\item Q4: V4 internal labels clean; integration must rewrite
  Lemma~5d and Theorem~1's Layer~2 proof to drop legacy $\beta'$.
\item Q5: Still missing — $\alpha$ explicit clause if multiplicity
  varies; channel multiplicity (C5.MULT); martingale-difference /
  i.i.d. increment structure (Hoeffding requires more than bounded
  range); LLR finiteness/support overlap.
\end{itemize}

\textbf{Changes from v4 (Round D fixes):}

\begin{itemize}
\item \textbf{Q1 fix (rigorous Hoeffding inversion).}  Replaced
  $R^2\log(1/\beta)/(2\deltastar^2)$ with
  $2R^2\log(1/\beta)/\deltastar^2$ (factor of~4).  Step~5 of the
  proof now uses the rigorous chain $T \geq 2A/\deltastar
  \Rightarrow (T\delta - A)^2 \geq T^2\delta^2/4
  \Rightarrow$ exponent $\geq T\delta^2/(2R^2)$, eliminating the
  asymptotic-approximation gap in v4.
\item \textbf{Q2 fix.  (C5.HOEFF) strengthened, (C5.MULT) added.}
  Replaced ``family-wise corrected'' with explicit fixed-cardinality
  monitored partition + uniform LLR clip.  Channel-multiplicity
  policy promoted to its own clause (C5.MULT) for clarity.
\item \textbf{Q3 fix.  Calibration-uncertainty clause $\bcal$.}
  Added explicit treatment: $\Ncasc \geq \Tbeta$ is certified at
  audit time as a hard constraint; calibration error is folded
  into a separate $\bcal$ probability if the operator chooses
  empirical estimation.  Total Layer~2 failure becomes
  $\beta + \bclk + \bcal$ in the most general form, $\beta + \bclk$
  if calibration is certified.
\item \textbf{Q4 fix.  Integration prerequisite list.}  Spelled out
  the Lemma~5d / main\_theorem.tex / deployment\_tooling.tex
  rewrites required at integration: drop $\beta'$ everywhere, use
  $\beta + \bclk + \bcal$ for total Layer~2 failure.
\item \textbf{Q5 fix.  Additional clauses.}
  \begin{itemize}
  \item (C5.IID): SPRT log-likelihood-ratio increments are
    martingale-differences under $p_1$ (or i.i.d. under fixed
    distribution), with bounded conditional variance proxy.
  \item (C5.SUPP): per-channel likelihood ratio satisfies
    $p_0(x), p_1(x) \geq \epsilon > 0$ for all $x$ in the
    monitored partition (or LLR is clipped to $[-R, R]$).
  \item $\alpha$ explicit: monitor Type-I rate, deployment-class
    if (C5.MULT) bounds multiplicity.
  \end{itemize}
\item \textbf{Lemma 4 prose nit.}  Round~D flagged that
  appendix's prose ``$\Lambda_t < A$ is sufficient'' should read
  ``necessary''.  This is an integration item for §A.4, not a
  Lemma~6 issue.
\end{itemize}

\textbf{Carried forward from v4 (unchanged in v5):}
\begin{itemize}
\item Supremum-form $h_{\mathrm{detect}}(\theta) := \sup_s
  \epsgap(s)$.
\item Discrete boundary convention $t_0 := t_0^-$.
\item $\rhogap$ rename + (C11) framing.
\item Union-class KL floor $\deltastar$.
\item (C11.CLK) deterministic floor + $\bclk$ structure.
\end{itemize}

\textbf{Round C other issues.}
\begin{itemize}
\item (C11.CLK) too weak: existence of $N_{\min}, q_{\min} > 0$
  doesn't imply $\Tbeta \leq \tau_{\mathrm{meta}}$ with calibrated
  probability.
\item $\beta'$ mislabeled: v3 used $\beta' = \exp(-\kappa
  \tau_{\mathrm{meta}} \deltastar)$ (the SPRT detection-tail form)
  for the cascade-clock event $\{\Tbeta \leq \tau_{\mathrm{meta}}\}$,
  conflating two distinct probabilities.
\item Constant taxonomy missing: $\alpha$ (Type-I rate), $R$ (LLR
  range), channel multiplicity policy, separate $\bclk$ for clock
  failure, separate from SPRT detection failure.
\item Q2 boundary convention sufficient inside Lemma 6 but must
  surface into Theorem~1's Layer~2 proof at integration time.
\end{itemize}

\textbf{Changes from v3 (Round C fixes):}

\begin{itemize}
\item \textbf{Q1 fix (algebra).}  Define $\Tbeta$ directly from
  Lemma~4's Hoeffding form:
  \[
    \Tbeta(\beta, \deltastar) \;:=\;
    \max\!\left\{\frac{2A}{\deltastar},\;
    \frac{R^2 \log(1/\beta)}{2 \deltastar^2}\right\},
  \]
  where $A := \log((1-\beta)/\alpha)$ is Paper~5's SPRT threshold
  and $R$ is the LLR-increment range from (C5.HOEFF).  In the
  asymptotic regime $\Tbeta \gg A/\deltastar$, the second term
  dominates and gives $\Pr[T_{\mathrm{detect}} > \Tbeta] \leq
  \beta$.  Critical: $\Tbeta$ scales as $1/\deltastar^2$, not
  $1/\deltastar$.
\item \textbf{(C5.HOEFF) added.}  Bounded LLR-increment range
  clause: per-step SPRT log-likelihood-ratio increments are bounded
  in absolute value by $R$, with $R$ deployment-class (not scaling
  with $|P|$ via channel multiplicity).
\item \textbf{(C11.CLK) strengthened.}  Replaced existence-only
  formulation with: there exists deterministic $\Ncasc$ such that
  $\Pr[\Nev(\tau_{\mathrm{meta}}) \geq \Ncasc] \geq 1 - \bclk$, and
  $\Ncasc \geq \Tbeta$ as an audit-time calibration constraint.
  This makes the bound $\Tbeta \leq \tau_{\mathrm{meta}}$
  operational with named failure probability $\bclk$.
\item \textbf{$\beta'$ relabeled.}  V3's use of $\beta' = \exp(-\kappa
  \tau_{\mathrm{meta}} \deltastar)$ is dropped from this lemma's
  statement.  The cascade-clock event $\{\Tbeta \leq
  \tau_{\mathrm{meta}}\}$ now has its own probability bound via
  $\bclk$ from (C11.CLK).  Total Layer~2 failure probability is
  $\beta + \bclk + \bcal$ (with $\bcal = 0$ in the certified case).
\item \textbf{Constant taxonomy completed.}  $\alpha$ surfaced as
  monitor-design parameter; $R$ surfaced via (C5.HOEFF); $\bclk$
  surfaced via (C11.CLK); channel multiplicity flagged as
  deployment-class via (C5.HOEFF) range constancy.
\item \textbf{Q2 integration note.}  V3's $t_0 := t_0^-$ convention
  carried forward; integration must update Theorem~1's Layer~2
  proof to express the Layer~1$\to$Layer~2 transition in SPRT
  exposure clock.
\end{itemize}

\textbf{Carried forward from v3 (unchanged in v4):}
\begin{itemize}
\item Supremum-form $h_{\mathrm{detect}}(\theta) := \sup_s
  \epsgap(s)$ matching Theorem~1's Layer~2 application.
\item Discrete boundary convention $t_0 := t_0^-$.
\item $\rhogap$ rename and reinforcement that $\rhogap$ excludes
  $K_{\mathrm{Lip}}$.
\item Union-class KL floor $\deltastar = \min(\deltaadv,
  \delta_{\mathrm{adv}}^{\mathrm{env}})$.
\item (C11) framed for promotion to theorem-level condition;
  source-citation correction re Paper~9.
\end{itemize}

\section{Setup: the monitoring clock and $h_{\mathrm{detect}}$}

Theorem~\ref{thm:deployment_safety}'s Layer~2 covers adversarial
strategies $q \in \Aadv \cup \Aadv^{\mathrm{env}}$ during the
detection window
$[t_0, t_0 + T_{\mathrm{detect}}]$.

\paragraph{Boundary convention.}  We define $t_0$ as the
\emph{last SPRT exposure step} at which the deployment is still
within Layer~1's safe region; equivalently, $t_0$ is the SPRT step
immediately preceding the first crossing event ($t_0 := t_0^-$ in
the v2 notation).  $T_{\mathrm{detect}}$ is the number of SPRT
steps from $t_0$ until the SPRT log-likelihood ratio crosses the
detection threshold (Lemma~\ref{lem:4}); the crossing step itself
is counted within $T_{\mathrm{detect}}$.

Under this convention, at $s = s_0 = t_0$ the Layer~1 bound applies
exactly: $\epsgap(t_0) \leq h_{\mathrm{static}}(\theta)$.  The
single-step jump that first crosses Layer~1 occurs at $s_1$, and is
bounded by $\rhogap(\deltastar)$ via (C11).  No ``absorbed buffer''
is required.

\paragraph{Detection-window bound.}
Layer~2 applies $h_{\mathrm{detect}}$ as a point-in-time bound:
\[
  |g(\Ttruth) - g(\Pproxy)|(s) \;\leq\; \mathrm{Lip}(g) \cdot
  h_{\mathrm{detect}}(\theta)
  \quad \text{for all } s \in [t_0, t_0 + T_{\mathrm{detect}}].
\]
Therefore $h_{\mathrm{detect}}$ must bound $\epsgap(s)$ at every
$s$ in the window:
\[
  h_{\mathrm{detect}}(\theta) \;:=\; \sup_{s \in
  [t_0, t_0 + T_{\mathrm{detect}}]} \epsgap(s).
\]

\paragraph{The monitoring clock and (C11.CLK) clock comparability.}
The SPRT machinery (Paper~5) operates on a per-event clock: each
ledger entry triggers an SPRT update, and the per-step KL divergence
drives detection.  All times in this lemma — $t_0$,
$T_{\mathrm{detect}}$, $\Tbeta$, and the cascade-time references
$\tau_{\mathrm{meta}}$, $T_{\mathrm{cascade}}$ — are expressed in
this SPRT exposure clock.

\begin{assumption}[(C5.SPRT): SPRT applicability]
\label{ass:C5SPRT}
The deployment's SPRT monitor satisfies the following four
sub-clauses, ensuring Lemma~\ref{lem:4}'s Hoeffding-Azuma step
applies with deployment-class constants:

\textbf{(C5.HOEFF) Clipped LLR with bounded Hoeffding range.}  The
raw per-step log-likelihood-ratio increments
$\ell_n^{\mathrm{raw}} := \log(p_1(X_n)/p_0(X_n))$ are clipped to
the symmetric interval $[-\Bclip, \Bclip]$:
\[
  \ell_n \;:=\; \mathrm{clip}\bigl(\ell_n^{\mathrm{raw}}, -\Bclip,
  \Bclip\bigr) \in [-\Bclip, \Bclip],
\]
with $\Bclip$ a deployment-class constant independent of $|P|$.
The Hoeffding range width used in Lemma~\ref{lem:4} is
\[
  \Rh \;:=\; 2\Bclip,
\]
matching Lemma~4's $R = b - a$ convention (interval width).
Clipping ensures the Hoeffding step applies regardless of
underlying-distribution support tails.

\textbf{(C5.MULT) Bounded channel multiplicity, fixed before
deployment.}  The deployment's channel-monitoring policy uses a
fixed-cardinality monitored partition: $\Kch \in \mathbb{N}$ is a
deployment-class constant independent of $|P|$, with familywise
$\alpha/\beta$ allocation over at most $\Kch$ channels.  The
monitored partition and channel-selection / testing policy are
fixed \emph{before} deployment; adaptive creation of new monitored
channels during deployment violates (C5.MULT).  In particular, the
four-channel agent-side observable (Paper~5) plus the
environment-side mirror (Paper~10 Lemma~5e) gives $\Kch \leq 8$.

\textbf{(C5.IID) Adapted increments with conditional drift floor.}
Under the alternative hypothesis $p_1$, the clipped per-step LLR
increments $\ell_n$ are adapted to filtration
$\{\mathcal{F}_n\}$ with conditional mean satisfying
\[
  \mathbb{E}[\ell_n \mid \mathcal{F}_{n-1}] \;\geq\; \delta_n
  \;\geq\; \deltastar \;>\; 0.
\]
The centered increments $\tilde\ell_n := \ell_n -
\mathbb{E}[\ell_n \mid \mathcal{F}_{n-1}]$ form a
martingale-difference sequence (bounded by $\Rh$ in range) to
which Azuma--Hoeffding applies.  The i.i.d.\ case (Paper~5
default) is the special case where conditional law is constant.

\textbf{(C5.SUPP) LLR finiteness.}  Almost-sure finiteness of
$\ell_n$ holds automatically under (C5.HOEFF) clipping.  This
sub-clause is a consistency check, not a new constraint.

\textbf{Important: $\deltastar$ is the post-clipping drift floor.}
Clipping can reduce drift from the raw KL value
($\delta_{\mathrm{adv}}, \delta_{\mathrm{adv}}^{\mathrm{env}}$
from Lemmas~5b/5e).  Operators must verify that the post-clipping
conditional-mean floor still satisfies $\mathbb{E}[\ell_n \mid
\mathcal{F}_{n-1}] \geq \deltastar$ for the clipped LLR; if not,
$\deltastar$ in this lemma should be replaced with the post-clipping
drift floor.
\end{assumption}

\begin{assumption}[(C11.CLK): Clock comparability with calibrated probability]
\label{ass:C11CLK}
There exist deployment-class constants $\Ncasc \in \mathbb{N}$ and
$\bclk \in (0,1)$ such that the SPRT exposure event-counting process
$\Nev(\cdot)$ satisfies:
\[
  \Pr\bigl[\Nev(\tau_{\mathrm{meta}}) \geq \Ncasc\bigr] \;\geq\;
  1 - \bclk,
\]
where $\Ncasc$ is certified at audit time to be a \emph{lower bound}
on the true event-throughput floor before cascade.  The operator
further certifies the audit-time inequality
\[
  \Ncasc \;\geq\; \Tbeta(\beta, \deltastar)
\]
holds (with $\Tbeta$ the SPRT detection quantile defined below), so
$\Tbeta \leq \Ncasc \leq \Nev(\tau_{\mathrm{meta}})$ with probability
$\geq 1 - \bclk$.

\textbf{Optional empirical-calibration variant.}  If the operator
estimates any of $\Tbeta$'s inputs ($\Bclip$ upper bound,
$\deltastar$ lower bound, $\alpha/\beta$ allocation) or $\Ncasc$
empirically rather than from conservative certificates, then a
calibration-uncertainty probability $\bcal \in (0,1)$ must satisfy:
\[
  \Pr\!\left[
  \begin{array}{l}
  \Bclip < \Bclip^{\mathrm{true}} \;\text{or}\;
  \deltastar > (\delta_*^{\mathrm{clip}})^{\mathrm{true}}\;\text{or}\\
  \Ncasc > \Ncasc^{\mathrm{true}}
  \end{array}
  \right] \;\leq\; \bcal.
\]
The directionality is uniform: each certificate must be a
\emph{conservative} (worst-case) bound — $\Bclip$ \emph{upper},
$\deltastar$ \emph{lower}, $\Ncasc$ \emph{lower}.  Total Layer~2
failure inflates to $\beta + \bclk + \bcal$.  In the certified-
conservative case (all inputs certified hard, no empirical
estimation), $\bcal = 0$.
\end{assumption}

\textbf{Why (C5.SPRT) and (C11.CLK) are required.}  Lemma~4's
Hoeffding-Azuma step requires:
\begin{itemize}
\item Bounded range $\Rh = 2\Bclip$ via (C5.HOEFF) clipping.
\item Constant cardinality $\Kch$ via (C5.MULT).
\item Adapted increments with drift floor $\deltastar$ via (C5.IID),
  measured \emph{post-clipping}.
\item Almost-sure finiteness via (C5.SUPP) — automatic under
  clipping.
\end{itemize}
(C11.CLK) maps detection time (in SPRT exposure clock) to wall-clock
cascade time via the event-throughput floor $\Ncasc$.  The
sub-clause bundling under (C5.SPRT) matches the (C9.*) and (C10.*)
patterns.

\textbf{Sign-direction discipline.}  All operator certificates point
in the safe direction: $\Bclip$ certified \emph{upper} (so clipped
LLR range is at most $\Rh$, making Hoeffding bound tighter is a gift
not a hazard), $\deltastar$ certified \emph{lower} (so detection
drift is at least $\deltastar$, making $\Tbeta$ bound conservative),
$\Ncasc$ certified \emph{lower} (so event throughput is at least
$\Ncasc$, making $\Tbeta \leq \Ncasc$ inequality robust).  V5 had
the empirical $\Ncasc$ direction backwards; v6 corrects this.

\section{Bounded gap-growth-rate condition (C11)}

\begin{assumption}[(C11): Bounded gap-growth rate per SPRT step]
\label{ass:C11}
For any adversarial strategy $q \in \Aadv \cup
\Aadv^{\mathrm{env}}$ producing channel deviations bounded by
$\deltastar$, the post-clipping union-class drift floor satisfying
$\deltastar \leq \min(\deltaadv,
\delta_{\mathrm{adv}}^{\mathrm{env}})$ — where $(\deltaadv,
\delta_{\mathrm{adv}}^{\mathrm{env}})$ are the raw Lemmas~5b/5e
KL floors and the inequality accounts for any drift reduction from
(C5.HOEFF) clipping — the per-SPRT-step proxy-truth gap growth is
bounded:
\[
  (\Delta \epsgap)_n \;:=\; \epsgap(s_n) - \epsgap(s_{n-1}) \;\leq\;
  \rhogap(\deltastar),
\]
where $\rhogap(\deltastar)$ is a deployment-class constant
depending on $\mathcal{D}$'s operational parameters (channel
sensitivity-to-gap mapping, event-rate normalization) but
independent of $|P|$.

\textbf{Continuous-time variant.}  Equivalently, in continuous
SPRT exposure clock $s$:
\[
  \frac{d \epsgap}{ds} \;\leq\; \rhogap(\deltastar).
\]
The discrete and continuous forms agree under fixed step duration
or bounded event-rate normalization.
\end{assumption}

\textbf{Operational interpretation.}  Channel deviations bounded by
$\deltastar$ (else SPRT fires faster) limit how much $\epsgap$ can
grow per SPRT step.  The constant $\rhogap$ characterizes the
deployment's channel-deviation-to-gap-growth coupling.

\textbf{$\rhogap$ is calibrated in the $\epsgap$ metric only.}
$\rhogap$ is a bound on the per-step proxy-truth gap process
$\epsgap = $ (Microfoundation's relative sup-norm), \emph{before}
applying any alignment-property mapping $g$.  In particular,
$\rhogap$ excludes the Lipschitz constant $K_{\mathrm{Lip}}$ of $g$;
the Layer~2 operational form multiplies $h_{\mathrm{detect}}$ by
$K_{\mathrm{Lip}}$ separately via (C6).  This eliminates any
double-counting between (C11) gap-dynamics and (C6) alignment-side
smoothness.

\textbf{(C11) is new content.}  \citep[Microfoundation]{lasser2026micro}'s
``$T$ failure modes'' remark notes adequacy/failure modes of the
operational truth $T$ but does not establish a formal per-step
gap-growth-rate bound.  (C11) is a Paper~10 operational assumption,
analogous to (C7) bounded co-evolution and SA1 HHI surrogate
adequacy: deployment-class condition with empirical testability,
not a derived consequence of source-paper machinery.

\section{Statement of Lemma 6 (v6)}

\begin{lemma}[$h_{\mathrm{detect}}$ intensivity]
\label{lem:6}
Under Theorem~\ref{thm:deployment_safety}'s conditions (C5)
continuous SPRT monitoring, (C5.SPRT) SPRT applicability (with
sub-clauses C5.HOEFF/MULT/IID/SUPP), (C11) bounded gap-growth
rate, and (C11.CLK) clock comparability with calibrated
probability, combined with Lemma~\ref{lem:4} (SPRT Wald--Hoeffding
tail bound) and Lemma~\ref{lem:1} (Layer~1 intensivity):

\textbf{Detection quantile.}  Define the SPRT high-probability
detection quantile via the conservative inversion of Lemma~4's
Hoeffding bound:
\[
  \Tbeta(\beta, \deltastar) \;:=\;
  \max\!\left\{\frac{2A}{\deltastar},\;
  \frac{2 \Rh^2 \log(1/\beta)}{\deltastar^2}\right\},
\]
where:
\begin{itemize}
\item $A := \log((1-\beta)/\alpha)$ is Paper~5's SPRT threshold;
\item $\alpha$ is the Type-I error rate (deployment-class monitor
  parameter; familywise allocated over $\Kch$ channels);
\item $\Rh = 2\Bclip$ is the Hoeffding range width from
  (C5.HOEFF) (Lemma~4's $R = b - a$ convention);
\item $\beta$ is the target detection-tail failure probability;
\item $\deltastar$ is the post-clipping union-class drift floor
  ((C5.IID), Lemmas~\ref{lem:5b}, \ref{lem:5e}).
\end{itemize}

(i) [Tail bound on $T_{\mathrm{detect}}$.]  $\Tbeta$ rigorously
satisfies $\Pr[T_{\mathrm{detect}} > \Tbeta] \leq \beta$ via
Lemma~\ref{lem:4}'s Hoeffding form (proved in Step~5).

(ii) [Detection-window supremum.]  On the event
$\{T_{\mathrm{detect}} \leq \Tbeta\}$ (probability at least
$1-\beta$):
\[
  h_{\mathrm{detect}}(\theta) \;:=\;
  \sup_{s \in [t_0, t_0 + T_{\mathrm{detect}}]} \epsgap(s)
  \;\leq\; h_{\mathrm{static}}(\theta) + \rhogap(\deltastar) \cdot
  \Tbeta.
\]

(iii) [Lead-time-before-cascade.]  By (C11.CLK)'s deterministic
floor $\Ncasc \geq \Tbeta$ and clock-failure probability $\bclk$
(plus calibration-uncertainty $\bcal$ if applicable):
\[
  \Pr\bigl[\Tbeta \leq \Nev(\tau_{\mathrm{meta}})\bigr] \;\geq\; 1 -
  \bclk - \bcal,
\]
i.e., $\Tbeta \leq \tau_{\mathrm{meta}}$ in SPRT exposure clock
with probability at least $1 - \bclk - \bcal$.  In the certified
case ($\bcal = 0$), this reduces to $1 - \bclk$.

(iv) [Layer~2 total failure bound.]  Total Layer~2 failure
probability (detection failure or cascade-clock miscalibration):
\[
  \Pr[\text{Layer~2 fails}] \;\leq\; \beta + \bclk + \bcal.
\]

\textbf{Intensivity.}  $h_{\mathrm{detect}}(\theta)$ is intensive
in $|P|$ over the deployment class $\mathcal{D}$:
\begin{itemize}
\item $h_{\mathrm{static}}(\theta)$: intensive by Lemma~\ref{lem:1}.
\item $\rhogap(\deltastar)$: intensive by (C11).
\item $\Tbeta = \max\{2A/\deltastar, 2\Rh^2\log(1/\beta)/\deltastar^2\}$:
  intensive because $A$ depends on $\alpha, \beta$ (monitor-design
  parameters; $\alpha$ deployment-class via (C5.MULT) familywise
  allocation), $\Rh = 2\Bclip$ is deployment-class via (C5.HOEFF)
  clipping, and $\deltastar$ is deployment-class (post-clipping
  drift floor via (C5.IID), Lemmas~5b/5e).
\item Total bound $h_{\mathrm{static}} + \rhogap \cdot \Tbeta$ is
  a sum of intensive quantities.
\end{itemize}
\end{lemma}

\section{Proof}

\begin{proof}
We bound $h_{\mathrm{detect}}$ as the supremum of $\epsgap$ over
the detection window, derive $\Tbeta$ directly from Lemma~4's
Hoeffding form, and apply (C11.CLK) to bound the cascade-clock
event.

\emph{Step 1 (per-step decomposition).}  Under the boundary
convention from §2, set $s_0 := t_0$ (the last safe SPRT step) and
$s_n := t_0 + n$ in SPRT exposure clock.  For any
$0 \leq n \leq T_{\mathrm{detect}}$:
\[
  \epsgap(s_n) \;=\; \epsgap(s_0) + \sum_{i=1}^{n}
  (\Delta \epsgap)_i.
\]

\emph{Step 2 (apply (C11)).}  By Assumption~\ref{ass:C11}, each
per-step increment satisfies $(\Delta \epsgap)_n \leq
\rhogap(\deltastar)$.  Therefore:
\[
  \epsgap(s_n) \;\leq\; \epsgap(s_0) + n \cdot
  \rhogap(\deltastar).
\]

\emph{Step 3 (boundary condition).}  By the boundary convention,
$s_0 = t_0$ is the last SPRT step at which the deployment is still
within Layer~1's safe region.  Layer~1's intensivity result
(Lemma~\ref{lem:1}) gives:
\[
  \epsgap(s_0) \;=\; \epsgap(t_0) \;\leq\;
  h_{\mathrm{static}}(\theta).
\]
The first crossing event occurs at $s_1$, where (C11) bounds the
single-step jump by $\rhogap(\deltastar)$.

\emph{Step 4 (supremum over window).}  Combining Steps~1--3:
\[
  h_{\mathrm{detect}}(\theta) \;=\;
  \sup_{0 \leq n \leq T_{\mathrm{detect}}} \epsgap(s_n)
  \;\leq\; h_{\mathrm{static}}(\theta) +
  T_{\mathrm{detect}} \cdot \rhogap(\deltastar).
\]

\emph{Step 5 (rigorous Hoeffding inversion to derive $\Tbeta$).}
By Lemma~\ref{lem:4} (\texttt{appendices/proofs.tex} §A.4), the
SPRT detection time satisfies the Hoeffding form for the
\emph{post-clipping} process with range $\Rh = 2\Bclip$:
\[
  \Pr[T_{\mathrm{detect}} > t] \;\leq\;
  \exp\!\left(-\frac{2(t \deltastar - A)^2}{t \Rh^2}\right)
  \quad \text{for } t > A/\deltastar.
\]
We derive $\Tbeta$ such that this bound is $\leq \beta$, using a
rigorous (not asymptotic) chain.

\emph{Substep 5a (regime condition).}  Suppose
$T \geq 2A/\deltastar$.  Then $T\deltastar - A \geq T\deltastar/2$,
so $T\deltastar - A > 0$ and the Hoeffding bound applies.

\emph{Substep 5b (squared lower bound).}  Squaring the inequality
from~5a:
\[
  (T\deltastar - A)^2 \;\geq\; \frac{T^2 \deltastar^2}{4}.
\]

\emph{Substep 5c (exponent lower bound).}  Substituting into the
Hoeffding exponent:
\[
  \frac{2(T\deltastar - A)^2}{T \Rh^2} \;\geq\;
  \frac{2 \cdot T^2 \deltastar^2 / 4}{T \Rh^2} \;=\;
  \frac{T \deltastar^2}{2 \Rh^2}.
\]

\emph{Substep 5d (solve for $T$ such that exponent $\geq \log(1/\beta)$).}
We want
$\frac{T \deltastar^2}{2 \Rh^2} \;\geq\; \log(1/\beta)$,
which gives
\[
  T \;\geq\; \frac{2 \Rh^2 \log(1/\beta)}{\deltastar^2}.
\]

\emph{Substep 5e (conservative $\Tbeta$).}  Combining the regime
condition (Substep~5a) and the tail condition (Substep~5d):
\[
  \Tbeta \;:=\; \max\!\left\{\frac{2A}{\deltastar},\;
  \frac{2\Rh^2 \log(1/\beta)}{\deltastar^2}\right\}.
\]
For any $T \geq \Tbeta$, both the regime condition and the
exponent-magnitude condition hold, so $\Pr[T_{\mathrm{detect}} > T]
\leq \exp(-T\deltastar^2/(2\Rh^2)) \leq \beta$.  In particular,
$\Pr[T_{\mathrm{detect}} > \Tbeta] \leq \beta$.

\emph{Range-width discipline ($R$ vs.\ $\Bclip$ vs.\ $\Rh$).}  V5
clipped to $[-R, R]$ (interval width $2R$) but used the symbol $R$
for the Hoeffding range — confusing range width with clip radius.
V6 uses $\Bclip$ for the (symmetric) clip radius and
$\Rh = 2\Bclip$ for the Hoeffding range, matching Lemma~4's
$R = b - a$ convention.  Concretely:
\[
  \Tbeta \;=\; \max\!\left\{\frac{2A}{\deltastar},\;
  \frac{2 \cdot (2\Bclip)^2 \log(1/\beta)}{\deltastar^2}\right\}
  \;=\; \max\!\left\{\frac{2A}{\deltastar},\;
  \frac{8 \Bclip^2 \log(1/\beta)}{\deltastar^2}\right\}.
\]

\emph{Cumulative factor tracking (v3 $\to$ v4 $\to$ v5 $\to$ v6).}
\begin{itemize}
\item v3: $\Tbeta \propto 1/\deltastar$ — wrong scaling.
\item v4: $\Tbeta \propto 1/(2\deltastar^2)$ — correct scaling but
  $\sim$$4\times$ too small (asymptotic approximation).
\item v5: $\Tbeta \propto 2/\deltastar^2$ — correct scaling and
  rigorous chain, but used $R$ ambiguously.
\item v6: $\Tbeta \propto 8\Bclip^2/\deltastar^2$ — clip radius
  $\Bclip$ and Hoeffding range $\Rh = 2\Bclip$ disambiguated;
  matches Lemma~4 convention exactly.
\end{itemize}

\emph{Step 6 (detection-window supremum bound).}  On the event
$\{T_{\mathrm{detect}} \leq \Tbeta\}$ (probability at least
$1-\beta$):
\[
  h_{\mathrm{detect}}(\theta) \;\leq\;
  h_{\mathrm{static}}(\theta) + \rhogap(\deltastar) \cdot
  \Tbeta.
\]
This proves part~(ii) of the lemma.

\emph{Step 7 (cascade-clock event via (C11.CLK)).}  Separately
from the SPRT tail, (C11.CLK) provides:
$\Pr[\Nev(\tau_{\mathrm{meta}}) \geq \Ncasc] \geq 1 - \bclk$, with
$\Ncasc$ certified at audit time as a \emph{lower} bound on the
true event-throughput floor ($\Ncasc \leq \Ncasc^{\mathrm{true}}$
deterministically in the certified case, with probability $\geq 1
- \bcal$ in the empirical case).  The audit-time inequality
$\Ncasc \geq \Tbeta$ is enforced as a hard constraint.

Combining:
\[
  \Pr[\Tbeta \leq \Nev(\tau_{\mathrm{meta}})] \;\geq\;
  \Pr[\Tbeta \leq \Ncasc \leq \Nev(\tau_{\mathrm{meta}})]
  \;\geq\; 1 - \bclk - \bcal.
\]
The first inequality uses the audit constraint $\Tbeta \leq
\Ncasc$ (deterministic).  The second uses (C11.CLK)'s event-count
floor (probability $\geq 1 - \bclk$) and the calibration
certificate (probability $\geq 1 - \bcal$, in the empirical case;
$\bcal = 0$ certified).  Union bound: $\bclk + \bcal$.

In the certified case ($\bcal = 0$), this reduces to $1 - \bclk$.
This proves part~(iii).

\emph{Sign-direction note.}  V5 stated
$\Pr[\Ncasc^{\mathrm{est}} \geq \Ncasc^{\mathrm{true}}] \geq 1 -
\bcal$, which is the wrong direction: an empirical estimate above
the true floor would not give a conservative bound.  V6 corrects
this by stating $\Ncasc$ (whether certified or empirical) is a
\emph{lower} bound on $\Ncasc^{\mathrm{true}}$, so
$\Ncasc \leq \Nev(\tau_{\mathrm{meta}})$ holds whenever the
event-throughput floor is met.

\emph{Step 8 (Layer 2 total failure bound).}  By union bound on
the failure modes:
\begin{itemize}
\item Detection-tail $\{T_{\mathrm{detect}} > \Tbeta\}$:
  probability $\leq \beta$ (Step~5).
\item Clock-failure $\{\Tbeta > \Nev(\tau_{\mathrm{meta}})\}$:
  probability $\leq \bclk + \bcal$ (Step~7).
\end{itemize}
\[
  \Pr[\text{Layer~2 fails}] \;\leq\; \beta + \bclk + \bcal.
\]
This proves part~(iv).

\emph{Step 9 (intensivity).}  Each constant is deployment-class
independent of $|P|$:
\begin{itemize}
\item $h_{\mathrm{static}}(\theta)$: intensive by Lemma~\ref{lem:1}
  (under (C6)--(C10)).
\item $\rhogap(\deltastar)$: deployment-class by (C11).
\item $\Tbeta = \max\{2A/\deltastar, 2\Rh^2\log(1/\beta)/\deltastar^2\}$:
  $A = \log((1-\beta)/\alpha)$ depends on $\alpha$ (deployment-class
  via (C5.MULT) familywise allocation over $\Kch$ channels) and
  $\beta$ (monitor-design parameter); $\Rh = 2\Bclip$ is
  deployment-class via (C5.HOEFF) clipping; $\deltastar$ is
  deployment-class via (C5.IID) post-clipping drift floor (rooted
  in Lemmas~5b/5e KL floors).  All inputs intensive, so $\Tbeta$
  intensive.
\item $\Ncasc, \bclk, \bcal$: deployment-class by (C11.CLK)
  (operator audit-time calibration).
\end{itemize}
Therefore $h_{\mathrm{detect}}(\theta) = h_{\mathrm{static}} +
\rhogap \cdot \Tbeta$ is intensive in $|P|$ over $\mathcal{D}$.
$\Box$
\end{proof}

\section{Discussion}

\paragraph{The two roles of cascade time (carried from v3).}
\begin{itemize}
\item \emph{Integration horizon for $\rhogap$:} $\Tbeta$, the SPRT
  high-probability detection quantile, gives a legitimate
  \emph{upper} bound on $T_{\mathrm{detect}}$ derived from
  Lemma~4's Hoeffding form.
\item \emph{Lead-time-before-cascade guarantee:}
  $\tau_{\mathrm{meta}} \leq T_{\mathrm{cascade}}$ (a \emph{lower}
  bound on cascade time, from
  \citep[Phase~Redundancy]{lasser2026phase}).  In v4, this is
  composed with $\Tbeta$ via (C11.CLK)'s deterministic floor
  $\Ncasc \geq \Tbeta$ and clock-failure probability $\bclk$ —
  \emph{not} via the SPRT detection-tail form (which was v3's
  mistake).
\end{itemize}

\paragraph{Why v3's $\beta'$ labeling was wrong.}
V3 labeled $\beta' = \exp(-\kappa \tau_{\mathrm{meta}} \deltastar)$
as the failure probability for the cascade-clock event $\{\Tbeta
\leq \tau_{\mathrm{meta}}\}$.  But this expression is the SPRT
\emph{detection-tail} form (probability that detection takes longer
than time $\tau_{\mathrm{meta}}$), not a cascade-clock
miscalibration probability.  V4 separates them: $\beta$ is the
detection-tail parameter; $\bclk$ is the clock-calibration failure
parameter (named in (C11.CLK)); their union bound (with $\bcal$
for the empirical-calibration variant) gives total Layer~2 failure
$\leq \beta + \bclk + \bcal$.

\paragraph{Why v3 and v4 had successive $\Tbeta$ errors.}
\emph{V3 error.}  V3 used $\Tbeta = \log(1/\beta)/(\kappa \delta)$
treating $\kappa$ as $\delta$-independent.  But Lemma~4's
$\kappa = 2\delta/R^2$ contains a factor of $\delta$, so the exponent
$\kappa T \delta = 2T\delta^2/R^2$ scales as $\delta^2$, not
$\delta$.  V4 corrected the scaling.

\emph{V4 error.}  V4 used $\Tbeta = R^2\log(1/\beta)/(2\deltastar^2)$
from a non-rigorous asymptotic approximation.  Substituting this
$T$ into the exact Hoeffding bound gives exponent $\log(1/\beta)
\cdot (1 - A\deltastar/c)^2$ with $c = R^2\log(1/\beta)/2$, strictly
\emph{less than} $\log(1/\beta)$ when $A > 0$.  V5 corrects this by
the rigorous chain in Step~5: $T \geq 2A/\deltastar \Rightarrow
(T\deltastar - A)^2 \geq T^2\deltastar^2/4 \Rightarrow$ exponent
$\geq T\deltastar^2/(2R^2)$.  Solving for exponent $\geq
\log(1/\beta)$ gives $T \geq 2R^2\log(1/\beta)/\deltastar^2$ — four
times v4's coefficient.

\paragraph{Lemma 5d integration item (Round B Q3 follow-up,
carried).}
At integration time, §4.10's Lemma~5d should be restated using the
union-class floor $\deltastar$.  This is a small amendment to §4.10
and §A.4 that should accompany Lemma~6 integration.

\paragraph{Constant taxonomy (v6 complete).}
\begin{itemize}
\item $h_{\mathrm{static}}(\theta)$: intensive composition
  (Lemma~\ref{lem:1}).
\item $\rhogap(\deltastar)$: \emph{Paper 10 operational constant}
  via (C11); calibrated by per-channel sensitivity-to-gap
  measurement (Audit~7).
\item $\Tbeta = \max\{2A/\deltastar, 2\Rh^2\log(1/\beta)/\deltastar^2\}$:
  rigorous SPRT-derived detection quantile.
\item $A = \log((1-\beta)/\alpha)$: Paper~5 SPRT threshold.
\item $\alpha$: Type-I error rate, monitor-design parameter
  (deployment-class via (C5.MULT) familywise allocation over $\Kch$
  channels).
\item $\beta$: Type-II error rate / detection-tail parameter,
  monitor-design parameter (deployment-class).
\item $\Bclip$: symmetric clip radius for LLR increments
  ($[-\Bclip, \Bclip]$), surfaced via (C5.HOEFF); deployment-class.
\item $\Rh = 2\Bclip$: Hoeffding range width (Lemma~4 convention
  $R = b - a$).
\item $\Kch$: monitored-channel cardinality, surfaced via
  (C5.MULT) (deployment-class, fixed before deployment).
\item $\deltastar$: post-clipping union-class drift floor, via
  (C5.IID) and Lemmas~5b/5e (deployment-class).
\item $\Ncasc$: deterministic event-count floor, surfaced via
  (C11.CLK) (audit-certified \emph{lower} bound on true throughput,
  with $\Ncasc \geq \Tbeta$ as hard audit constraint).
\item $\bclk$: clock-calibration failure probability ((C11.CLK)).
\item $\bcal$: calibration-uncertainty failure probability
  ((C11.CLK) optional empirical variant); $\bcal = 0$ in certified
  case.  Sign-direction discipline: each $\Tbeta$-input certificate
  is conservative ($\Bclip$ upper, $\deltastar$ lower, $\Ncasc$
  lower).
\item $\tau_{\mathrm{meta}}, T_{\mathrm{cascade}}$: source-derived
  cascade-time references from
  \citep[Phase~Redundancy]{lasser2026phase} (lower bound only;
  with the (C7) bounded co-evolution caveat).
\end{itemize}

\paragraph{Layer 2 operational form.}
Combining Lemma~6 with $\mathrm{Lip}(g) \leq K_{\mathrm{Lip}}$ from
(C6):
\[
  |g(\Ttruth) - g(\Pproxy)|(s) \;\leq\; K_{\mathrm{Lip}} \cdot
  \bigl(h_{\mathrm{static}}(\theta) + \rhogap(\deltastar) \cdot
  \Tbeta\bigr)
\]
with probability at least $1 - (\beta + \bclk + \bcal)$ over the
joint event $\{T_{\mathrm{detect}} \leq \Tbeta\} \cap \{\Tbeta \leq
\Nev(\tau_{\mathrm{meta}})\}$.  The right-hand side is the sum of
intensive constants times $K_{\mathrm{Lip}}$, all $|P|$-independent.

\paragraph{Integration prerequisites (v6 expanded).}
At appendix integration, the following must accompany the v6 proof:

\textbf{Notation and front-matter.}
\begin{itemize}
\item §2 notation table: add $\Tbeta, \Ncasc, \bclk, \bcal, \Bclip,
  \Rh, \Kch, \deltastar, A$.
\item Front-matter introduction: drop legacy $\beta'$ references.
\end{itemize}

\textbf{Lemmas / appendix.}
\begin{itemize}
\item Update §4.10 Lemma~5d \emph{statement} (not just appendix
  prose) to use $\beta + \bclk + \bcal$ union-class form.
\item Update §4 Lemma~4 statement to make $\Rh = b - a$ convention
  explicit.
\item §A.4 prose nit: ``$\Lambda_t < A$'' is necessary for
  $T_{\mathrm{detect}} > t$, not sufficient.
\item Update Theorem~1 Layer~2 proof in §5
  (\texttt{main\_theorem.tex}) to use $\Tbeta$ instead of
  $T_{\mathrm{cascade}}$ in Step~6, adopt $t_0 := t_0^-$
  convention, cite Lemma~6 parts~(ii)--(iv).
\item Drop v1 source-citation overstatement in
  \texttt{main\_theorem.tex} line~365 (Paper~9's ``$T$ failure
  modes'').
\end{itemize}

\textbf{Theorem 1 condition list.}
\begin{itemize}
\item Add (C5.SPRT) SPRT applicability with sub-clauses
  (C5.HOEFF / MULT / IID / SUPP), bundled like (C9.*) and (C10.*).
\item Add (C11) bounded gap-growth rate.
\item Add (C11.CLK) clock comparability with calibrated
  probability.
\end{itemize}

\textbf{§3 invariants connection.}
Connect (C11.CLK) to $I_{11}$ (time-asymmetry): the event-throughput
floor $\Ncasc$ is a quantitative refinement of $I_{11}$'s qualitative
asymmetry assertion.  Whether to make this a separate twelfth
invariant or a sub-clause of $I_{11}$ is an editorial choice.

\textbf{§8 Audit 7 expanded.}
Audit~7 should cover: $\rhogap(\deltastar)$ calibration; $\Bclip$
upper-bound certification (LLR clipping radius); $\Rh = 2\Bclip$
range-width convention check; $\Kch$ channel-cardinality fix-before-
deployment; familywise $\alpha/\beta$ allocation; clipped-process
$\deltastar$ post-clipping verification; $\Ncasc$ event-throughput
floor lower-bound certification; $\bclk$ clock-failure probability;
$\bcal$ calibration-uncertainty (if empirical inputs).

\textbf{§6 cooperative-anchoring.}
Layer~2 failure decomposition $\beta + \bclk + \bcal$ has no direct
cooperative-anchoring implications (cooperative anchoring concerns
the pairwise substrate cooperative floor in Layer~1, not the
SPRT-monitor side of Layer~2).  No §6 update needed.

\section{What this proof does and does not establish}

\textbf{Establishes:}
\begin{itemize}
\item Algebraically correct $\Tbeta$ derivation from Lemma~4's
  Hoeffding form, with explicit $1/\deltastar^2$ scaling.
\item (C5.HOEFF) bounded LLR-increment range as a deployment-class
  condition, ensuring $R$ does not scale with $|P|$ via channel
  multiplicity.
\item (C11.CLK) strengthened with deterministic floor $\Ncasc \geq
  \Tbeta$ and named clock-failure probability $\bclk$, separating
  cascade-clock event from SPRT detection-tail event.
\item Supremum-form $h_{\mathrm{detect}}$ definition matching
  Theorem 1's Layer 2 application.
\item Rigorous discrete-time boundary condition under the
  $t_0 := t_0^-$ convention.
\item Union-class KL floor $\deltastar$ for Layer 2's full
  detection class.
\item Intensive bound on $h_{\mathrm{detect}}$ as $h_{\mathrm{static}}
  + \rhogap \Tbeta$ with high probability $1-\beta$.
\item Layer~2 total failure probability $\leq \beta + \bclk + \bcal$
  via union bound (with $\bcal = 0$ in the certified case).
\item Constant taxonomy with all $|P|$-dependent constants surfaced
  ($\alpha, \beta, \Bclip, \Rh, \Kch, \Ncasc, \bclk, \bcal$ all
  deployment-class via named conditions).
\end{itemize}

\textbf{Does not establish (deferred):}
\begin{itemize}
\item Tightness of $\rhogap(\deltastar), R, \Ncasc$.  Operators
  may calibrate tighter bounds for restricted threat models.
\item Existence of deployment configurations satisfying (C5.HOEFF),
  (C11), and (C11.CLK).  Verified empirically (Audit~7).
\item Structural derivation of (C11) from finer source-paper
  machinery.  Currently a Paper 10 operational assumption.
\item Layer~2 proof updates in Theorem~1 (must adopt $t_0 := t_0^-$
  convention and use $\Tbeta$ instead of $T_{\mathrm{cascade}}$ in
  Step~6 of the Layer~2 proof).  This is an integration task.
\end{itemize}

\section{Specific verification questions for v6 codex review}

\begin{enumerate}
\item \textbf{(Q1 range-width disambiguation.)}  V6 separates clip
  radius $\Bclip$ (symbol for symmetric clip $[-\Bclip, \Bclip]$)
  from Hoeffding range $\Rh = 2\Bclip$ (Lemma~4's $R = b - a$
  convention).  Step~5 uses $\Rh$ throughout:
  $\Tbeta = \max\{2A/\deltastar, 2\Rh^2 \log(1/\beta) /
  \deltastar^2\} = \max\{2A/\deltastar, 8\Bclip^2 \log(1/\beta) /
  \deltastar^2\}$.

  Is the range-width discipline now consistent end-to-end?  In
  particular: does the appendix Lemma~4 statement (referenced as
  $R = b - a$) need to be updated to use the $\Rh$ symbol
  explicitly, or is the convention sufficiently anchored by the
  v6 prose?

\item \textbf{(Q2 (C5.SPRT) bundling and (C5.MULT) adversarial
  clause.)}  V6 bundles HOEFF/MULT/IID/SUPP under (C5.SPRT) with
  sub-clauses, matching (C9.*) and (C10.*).  (C5.MULT) now
  includes ``monitored partition fixed before deployment; adaptive
  channel creation falls into residual class.''

  Is the bundling clean, and is the adversarial-channel clause
  sufficient?  In particular: does adversarial \emph{selection
  among already-monitored channels} (not creation of new ones)
  need additional treatment, or is that already covered by the
  per-channel KL floor structure of Lemmas~5b/5e?

\item \textbf{(Q3 (C5.IID) reformulation.)}  V6 reformulates
  (C5.IID) to avoid v5's contradiction (martingale-difference +
  positive drift): now adapted increments with conditional mean
  $\mathbb{E}[\ell_n \mid \mathcal{F}_{n-1}] \geq \delta_n \geq
  \deltastar$, with the centered increments
  $\tilde\ell_n := \ell_n - \mathbb{E}[\ell_n \mid
  \mathcal{F}_{n-1}]$ being the martingale-difference sequence.
  Plus explicit note that $\deltastar$ is the post-clipping drift
  floor (clipping can reduce drift from raw KL).

  Is this the right framing for Hoeffding-Azuma to apply?  Does
  the post-clipping drift verification need its own audit
  (Audit~7 sub-item) or is it implicit in the (C5.SPRT) bundle?

\item \textbf{(Q4 empirical-calibration sign fix.)}  V6 corrects
  v5's sign error: $\Ncasc$ is now certified \emph{lower} bound
  on true throughput, $\Bclip$ is certified \emph{upper} bound on
  clip radius, $\deltastar$ is certified \emph{lower} bound on
  drift.  All certificates conservative with union failure
  probability $\bcal$.

  Is the sign-direction discipline now correct end-to-end?  Are
  there other places in the proof where sign-direction issues
  could lurk (e.g., in the post-clipping drift verification, or
  in the relationship between raw KL floor and clipped drift)?

\item \textbf{(Q5 integration-prerequisite expansion.)}  V6's
  ``Integration prerequisites'' paragraph now covers: §2 notation
  table, §4.10 Lemma~5d statement, §4 Lemma~4 statement (with
  $R = b - a$ convention), §A.4 prose nit, Theorem~1 Layer~2 proof
  rewrite, legacy $\beta'$ drop, (C5.SPRT)/(C11)/(C11.CLK)
  additions, §3 invariants/$I_{11}$ connection, expanded Audit~7,
  §6 explicitly noted as not needing update.

  Is the list now complete enough to drive the integration patch?
  Or are there integration items still missing?
\end{enumerate}

\end{document}
